% SPDX-License-Identifier: LPPL-1.3c
% Copyright (C) 2026 Christophe Jorssen
% Author and Current Maintainer: Christophe Jorssen <christophe.jorssen@gmail.com>
% LPPL maintenance status: maintained
% This file is part of luafemm. See LICENSE and LICENSES.md for its terms.
% Native TikZ keys, evaluated soft paths, and field/mesh pictures.
% Internal file: loaded by luafemm.tex with @ a letter.
% Physical paths are captured in a postaction, after PGF has resolved nodes,
% transformations and rounded corners. All dimensional conventions live in
% the manual; numerical values cross the Lua boundary as escaped strings.
\newif\ifluafemm@arrows\luafemm@arrowstrue
% Empty for ordinary regions; component winding paths set this locally.
\def\luafemm@componentturns{}
% Cache configuration belongs to the picture, before its begin hook runs.
% The choice handler stores a value rather than expanding user text as Lua.
\pgfkeys{/femm/cache/.is family,/femm/cache,
  file/.initial=\jobname-femm,mode/.initial=off,mode/.is choice,
  mode/off/.code={\pgfkeyssetvalue{/femm/cache/mode}{off}},
  mode/auto/.code={\pgfkeyssetvalue{/femm/cache/mode}{auto}},
  mode/refresh/.code={\pgfkeyssetvalue{/femm/cache/mode}{refresh}},
  mode/frozen/.code={\pgfkeyssetvalue{/femm/cache/mode}{frozen}}}
\pgfkeys{/femm/problem/.is family,/femm/problem,
  xmin/.initial=-90,xmax/.initial=90,ymin/.initial=-75,ymax/.initial=95,
  mesh size/.initial=2,mesher/.initial=delaunay,
  mesh tolerance/.initial=0,
  minimum angle/.initial=0,maximum area/.initial=0,max Steiner points/.initial=5000,
  tolerance/.initial=0.0000001,curve tolerance/.initial=.03,
  max iterations/.initial=60,
  /femm/region/.is family,/femm/region,
  material/.initial=air,library/.initial=,current density/.initial=material,curve tolerance/.initial=.03,
  magnetization angle/.initial=0,mesh size/.initial=0}
\def\luafemm@problemvalue#1{\pgfkeysvalueof{/femm/problem/#1}}
\def\luafemm@regionvalue#1{\pgfkeysvalueof{/femm/region/#1}}
\tikzset{
  femm/cache/.code={\pgfqkeys{/femm/cache}{mode=auto,#1}},
  femm/cache/.default={},
  femm/problem/.code={%
    \pgfqkeys{/femm/problem}{#1}%
    \tikzset{x=1mm,y=1mm,execute at begin picture={\luafemm@starttikz}}},
  femm/problem/.default={},
  femm/materials/.is family,
  femm/materials/mu r/.initial=,
  femm/materials/bh/.initial={},
  femm/materials/remanence/.initial=,
  femm/materials/coercivity/.initial=,
  femm/materials/library/.is choice,
  femm/materials/air/.style={mu r=1},
  femm/region/.code={%
    \pgfqkeys{/femm/region}{#1}%
    \tikzset{postaction={femm/capture}}},
  femm/region/.default={},
  femm/material/.code={\pgfkeys{/femm/region/material={#1}}},
  femm/current density/.code={\pgfkeyssetvalue{/femm/region/current density}{#1}},
  femm/magnetization angle/.code={\pgfkeyssetvalue{/femm/region/magnetization angle}{#1}},
  femm/mesh size/.code={\pgfkeyssetvalue{/femm/region/mesh size}{#1}},
  femm/curve tolerance/.code={\pgfkeyssetvalue{/femm/region/curve tolerance}{#1}},
  % This private postaction sees only the evaluated physical path, even
  % when the original path has nodes or additional drawing postactions.
  femm/capture/.code={\tikz@addmode{\luafemm@capturetikz}},
  femm/solve/.code={\tikz@addmode{\femmsolve}},
  femm/lines/.initial=29,
  femm/arrows/.is if=luafemm@arrows,
  femm/arrow spacing/.initial=22mm,
  femm/arrow length/.initial=1.3mm,
  femm/arrow width/.initial=.85mm,
  femm oriented field/.code={\ifluafemm@arrows
    % Nearly coincident contour vertices can overflow fixed-point veclen.
    % Keep the exact contour geometry and use PGF's floating length here.
    \tikzset{/pgf/fpu/install only={veclen},postaction={decorate},decoration={markings,
      mark=between positions .12 and .95 step \pgfkeysvalueof{/tikz/femm/arrow spacing}
        with {\arrow{Stealth[length=\pgfkeysvalueof{/tikz/femm/arrow length},
          width=\pgfkeysvalueof{/tikz/femm/arrow width}]}}}}\fi},
  femm field/.style={blue!65!black,line width=.32pt,femm oriented field},
  femm/field style/.code={\tikzset{femm current field/.style={#1}}},
  femm/mesh style/.code={\tikzset{femm current mesh/.style={#1}}},
  femm current field/.style={},femm current mesh/.style={},
  femm mesh/.style={draw=black!25,line width=.12pt},
  pics/femm field/.style={code={%
    \begingroup\pgftransformreset
    \directlua{luafemm.tex.tikz_field_picture(luafemm.current,
      luafemm.tex.number('\luaescapestring{\pgfkeysvalueof{/tikz/femm/lines}}','lines'))}%
    \endgroup}},
  pics/femm mesh/.style={code={%
    \begingroup\pgftransformreset
    \directlua{luafemm.tex.tikz_mesh_picture(luafemm.current)}%
    \endgroup}}
}
% Catalogue choices and existing user material styles share the selector.
\def\luafemm@selectlibrary#1{%
  \pgfkeyssetvalue{/femm/region/material}{#1}%
  \pgfkeyssetvalue{/femm/region/library}{#1}}
\pgfkeys{/femm/region/material/.is choice,
  /femm/region/material/.unknown/.code={%
    \edef\luafemm@selected{\pgfkeyscurrentname}%
    \pgfkeysifdefined{/tikz/femm/materials/\luafemm@selected/.@cmd}{%
      \pgfkeyslet{/femm/region/material}\luafemm@selected
      \pgfkeyssetvalue{/femm/region/library}{}%
    }{\pgfkeys@error{Unknown luafemm material '\luafemm@selected'. See docs/materials-catalog.md}}}}
\directlua{luafemm.tex.register_material_choices()}%
% Capture the picture frame only after TikZ has installed all its options.
\def\luafemm@starttikz{%
  % Several styles may set femm/problem. Their final values describe one
  % model; repeated begin hooks must not replace it or discard its boundaries.
  \ifluafemm@modelstarted\else
  \pgfgettransform\luafemm@roottransform
  \pgfkeyssetvalue{/femm/region/curve tolerance}{\luafemm@problemvalue{curve tolerance}}%
  \directlua{
    luafemm.tex.new_problem({
      xmin='\luaescapestring{\luafemm@problemvalue{xmin}}',xmax='\luaescapestring{\luafemm@problemvalue{xmax}}',
      ymin='\luaescapestring{\luafemm@problemvalue{ymin}}',ymax='\luaescapestring{\luafemm@problemvalue{ymax}}',
      h='\luaescapestring{\luafemm@problemvalue{mesh size}}',mesher='\luaescapestring{\luafemm@problemvalue{mesher}}',
      mesh_tolerance='\luaescapestring{\luafemm@problemvalue{mesh tolerance}}',
      min_angle='\luaescapestring{\luafemm@problemvalue{minimum angle}}',
      max_area='\luaescapestring{\luafemm@problemvalue{maximum area}}',
      max_steiner='\luaescapestring{\luafemm@problemvalue{max Steiner points}}',
      tolerance='\luaescapestring{\luafemm@problemvalue{tolerance}}',max_newton='\luaescapestring{\luafemm@problemvalue{max iterations}}'},
      '\luaescapestring{\luafemm@roottransform}',
      {'\the\pgf@xx','\the\pgf@xy','\the\pgf@yx','\the\pgf@yy'},
      {file='\luaescapestring{\pgfkeysvalueof{/femm/cache/file}}',
       mode='\luaescapestring{\pgfkeysvalueof{/femm/cache/mode}}'})}%
  \luafemm@modelstartedtrue
  \luafemm@pendingboundaries
  \let\luafemm@pendingboundaries\pgfutil@empty
  \fi
}
\def\luafemm@capturetikz{%
  \begingroup
  \pgfsyssoftpath@getcurrentpath\luafemm@capturedpath
  \pgfprocessround{\luafemm@capturedpath}{\luafemm@capturedpath}%
  \edef\luafemm@library{\luafemm@regionvalue{library}}%
  \pgfqkeys{/tikz/femm/materials}{mu r={},bh={},remanence={},coercivity={}}%
  \ifx\luafemm@library\pgfutil@empty
    \pgfqkeys{/tikz/femm/materials}{\luafemm@regionvalue{material}}%
  \fi
  \directlua{luafemm.tex.material('\luaescapestring{\luafemm@regionvalue{material}}',
    '\luaescapestring{\luafemm@library}',
    '\luaescapestring{\pgfkeysvalueof{/tikz/femm/materials/mu r}}',
    '\luaescapestring{\pgfkeysvalueof{/tikz/femm/materials/bh}}',
    '\luaescapestring{\pgfkeysvalueof{/tikz/femm/materials/remanence}}',
    '\luaescapestring{\pgfkeysvalueof{/tikz/femm/materials/coercivity}}')}%
  \edef\luafemm@currentinput{\luafemm@regionvalue{current density}}\def\luafemm@frommaterial{material}%
  \ifx\luafemm@currentinput\luafemm@frommaterial
    \edef\luafemm@current{\directlua{tex.sprint(luafemm.current.materials['\luaescapestring{\luafemm@regionvalue{material}}'].current)}}%
  \else
    \pgfkeys{/pgf/fpu=true}%
    \pgfmathparse{\luafemm@currentinput}%
    \pgfmathfloattofixed{\pgfmathresult}\let\luafemm@current\pgfmathresult
    \pgfkeys{/pgf/fpu=false}%
  \fi
  \pgfmathparse{\luafemm@regionvalue{magnetization angle}}\let\luafemm@angle\pgfmathresult
  \directlua{luafemm.tex.capture_path(luafemm.current,
    '\luaescapestring{\meaning\luafemm@capturedpath}',
    '\luaescapestring{\luafemm@regionvalue{material}}',
    luafemm.tex.number('\luaescapestring{\luafemm@current}','current density'),
    luafemm.tex.number('\luaescapestring{\luafemm@regionvalue{curve tolerance}}','curve tolerance'),
    luafemm.tex.number('\luaescapestring{\luafemm@angle}','magnetization angle'),
    luafemm.tex.number('\luaescapestring{\luafemm@regionvalue{mesh size}}','mesh size'),
    '\luaescapestring{\luafemm@componentturns}',
    '\ifpgfsys@eorule even odd\else nonzero\fi')}%
  \endgroup
}
\endinput
